\scanpage{98}
\mspage{25}
\diagram{\begin{tikzcd}[ampersand replacement=\&,column sep=large]
\begin{gathered}\text{Polygônes combinatoires }\widetilde P_{2n}\\
\text{d'ordre }2n,\text{ avec un}\\
\text{sous-polygône d'ordre}\\
n\text{ marqué, soit }P'_n
\end{gathered}\arrow[r,"\approx"]\&
\begin{gathered}\text{polygônes}\\\text{combinatoires }P_n\\\text{d'ordre }n\end{gathered}\\
(\widetilde P_{2n},P'_n)\arrow[r,mapsto]\&P'_n
\end{tikzcd}}
(9)

($n$ \uline{impair})

Donc la \deleted{représentation} \inserted{opération} naturelle de $D_{P_n}$ sur $\Sigma$, ou même sur $\boldsymbol\Sigma_{P_{2n}}$, \inserted{$=\mathbb P(V_{\widetilde P_{2n}})$,} se relève en une \deleted{représentation} \inserted{opération linéaire} de $D_{P_n}$ sur $V_{\widetilde P_{2n}}$. En fait, on \deleted{peut \unsure{même} préciser \unsure{plus}} \inserted{a}
\[
\tag{10}
V_{\widetilde P_{2n}}=\underbrace{V_{P'_n}\simeq V_{P_n}}_{\text{car }P'_n\simeq P_n}
\]
et l'action de $D_n$ sur $V_{\widetilde P_{2n}}$ peut s'identifier à son action sur $V_{P_n}$,
\deleted{et $\Sigma_{P_{2n}}$ — Ici $\boldsymbol\Sigma_{\widetilde P_{2n}}$ s'identifie à \unsure{une} la droite projective —}
\inserted{L'isomorphisme (10)} fournit d'ailleurs
\[
\tag{11}
\underset{\substack{\Vert\\\Sigma_{P_n}}}{\boldsymbol\Sigma_{\widetilde P_{2n}}}
=\mathbb P(V_{\widetilde P_{2n}})\overset\sim\longrightarrow\mathbb P(V_{P_n}),
\qquad(n\ \text{\uline{impair}})
\]
\deleted{ce qui est} et cette explicitation de $\Sigma_{P_n}$ comme $\mathbb P(V_{P_n})$ est \deleted{bien} effet évident a priori, car l'application
\[
\left\{
\begin{aligned}
S=S_{P_n}&\longrightarrow\mathbb P(V_{P_n})(\mathbb R)\\
s&\longmapsto\mathbb R s
\end{aligned}
\right.
\]
est \uline{injective} pour $n$ impair — de sorte que

\margin{97}
